\section{Related Work}
\label{sec:related}

\noindent\textbf{Static benchmarks.}
Text-to-SQL progress has been driven by single-turn benchmarks. Spider~\citep{yu2018spider} established cross-domain generalization; BIRD~\citep{li2023bird} scaled evaluation to messy, knowledge-grounded databases across professional domains; and Spider~2.0~\citep{lei2025spider2} pushes into enterprise workflows with multiple SQL dialects. These benchmarks advance query generation under the assumption of a well-formed question, leaving multi-turn interaction untested.

\namedpara{Conversational and multi-turn benchmarks.}
Related work brings dialogue into text-to-SQL evaluation. SParC~\citep{yu2019sparc} and CoSQL~\citep{yu2019cosql} introduce context-dependent question sequences and Wizard-of-Oz dialogues over Spider databases. More recent benchmarks raise interactive realism: BIRD-Interact~\citep{huo2026interact} couples each database with a function-driven user simulator and a CRUD-spanning task suite, and DySQL-Bench~\citep{sun2025dysql} evaluates dynamic refinement under evolving user intent. These benchmarks mainly measure end-to-end multi-turn execution, but none specifically targets the \emph{conditional dependency} between clarification layers or the \emph{non-monotonic retraction} of committed elements as evaluation dimensions. \textsc{TIDE-Bench} addresses these two phenomena directly, with a paired factorial design that attributes performance differences to chain ambiguity and user intent drift.

\namedpara{Ambiguity and clarification.}
Prior work treats ambiguity as central. AmbiQT~\citep{bhaskar2023ambiqt} characterizes lexical and structural ambiguity from overlapping schema names and join paths; AMBROSIA~\citep{saparina2024ambrosia} grounds three linguistic ambiguity types in controlled databases; PRACTIQ~\citep{dong2025practiq} pairs ambiguous and unanswerable questions with single-turn clarifications; and AmbiSQL~\citep{ding2026ambisql} couples a fine-grained taxonomy with interactive multiple-choice resolution. HiL-Bench~\citep{trinh2026hilbench} studies help-seeking across progressively discovered but independent blockers. These ambiguity-focused benchmarks define rich turn-local taxonomies for individual user utterances. \textsc{TIDE-Bench} is complementary: it studies how ambiguities depend on one another across turns (chain) and how committed intent can be retracted across turns (drift), two dialogue-level phenomena beyond these turn-local taxonomies.
